\documentclass[11pt]{article}
\usepackage{amsmath,amssymb,amsfonts}
\usepackage{graphicx}
\usepackage{hyperref}
\usepackage{booktabs}
\usepackage{geometry}
\geometry{margin=1in}

\title{Variational Generation of Relational Structure\\in the Relational Emergence Model:\\{\normalsize Quadratic Dynamical Cost and Regime Competition}}
\author{Yoshifumi Maruko\\Independent Researcher\\Yamagata, Japan}
\date{2026}

\begin{document}
\maketitle

\begin{abstract}
Relational Quantum Mechanics (RQM) defines physical states relationally but does not specify how relational structure itself becomes articulated. The Relational Emergence Model (REM) introduces a generative layer beneath relational description.

We model relational articulation as a variational selection problem over Hilbert-space factorizations. The selection functional takes the form
\[
\Phi = \Phi_S - \lambda C_H,
\qquad
C_H = \mathrm{Tr}(\rho H_{\partial F}^2) \geq 0,
\]
combining informational differentiation and a sign-definite boundary-fluctuation cost.

The parameter $\lambda\geq0$ acts as an external control parameter governing the balance between informational structure and dynamical stability. A three-qubit example demonstrates regime competition between informationally optimal and dynamically stabilized factorizations.
\end{abstract}

\section{Introduction}
Relational Quantum Mechanics proposes that physical states exist only relative to other systems. While this removes observer-independent states, it leaves open a prior question: how relational structure itself becomes articulated.

The Relational Emergence Model introduces a generative layer beneath relational description. Rather than assuming subsystem structure, REM models relational articulation as a structured selection problem. Relational structure is not assumed; it is selected.

\section{Variational Selection Principle}
Hilbert space admits multiple inequivalent tensor-product structures:
\[
\mathcal H \cong \mathcal H_A \otimes \mathcal H_B.
\]
Candidate factorizations form a space
\[
\mathcal F\cong U(N)/(U(n_A)\otimes U(n_B)).
\]
We define
\[
F^*=\arg\max_F\Phi(F),
\]
with
\[
\Phi=\Phi_S-\lambda C_H.
\]
The informational contribution is
\[
\Phi_S=I(A:B),
\]
and the dynamical cost is
\[
C_H=\mathrm{Tr}(\rho H_{\partial F}^2)\geq0.
\]
The equivalent minimization form is
\[
\mathcal L=-\Phi_S+\lambda C_H.
\]

\section{Information-Theoretic Quantities}
For $\rho_{AB}=|\Psi\rangle\langle\Psi|$,
\[
I(A:B)=S(\rho_A)+S(\rho_B)-S(\rho_{AB}),
\]
where
\[
S(\rho)=-\mathrm{Tr}(\rho\log_2\rho).
\]
For pure bipartitions,
\[
I(A:B)=2S(\rho_A).
\]

\section{Boundary Hamiltonian}
For $H=\sum_\alpha h_\alpha$, define
\[
H_{\partial F}=\sum_{\alpha:\,\mathrm{supp}(h_\alpha)\cap A\neq\varnothing,\,\mathrm{supp}(h_\alpha)\cap B\neq\varnothing} h_\alpha.
\]
Thus
\[
H=H_A+H_B+H_{\partial F}.
\]
The quadratic cost measures susceptibility to boundary-crossing fluctuations rather than mean boundary energy.

\section{Three-Qubit Illustration}
Consider
\[
|\Psi\rangle=\sqrt{0.6}|000\rangle+\sqrt{0.3}|111\rangle+\sqrt{0.1}|101\rangle.
\]
For candidate bipartitions $A|BC$, $B|AC$, and $C|AB$, the informational and dynamical terms may favor different factorizations. The selected structure changes where
\[
\Phi(F_1,\lambda^*)=\Phi(F_2,\lambda^*),
\]
so
\[
\lambda^*=\frac{\Phi_S(F_1)-\Phi_S(F_2)}{C_H(F_1)-C_H(F_2)}.
\]
The numerical value is model dependent and is not proposed as a universal constant.

\section{Physical Interpretation of $\lambda$}
$\lambda$ is an effective tradeoff parameter rather than a fundamental universal constant. Candidate interpretations include an inverse-temperature-like parameter, a ratio of decoherence and internal dynamical timescales, and a scale-dependent coarse-graining parameter.

A microscopic derivation from an explicit system--environment model remains future work.

\section{Relation to REM5}
REM5 replaces a sign-sensitive linear boundary expectation with a positive quadratic interaction cost. The form
\[
C_H=\mathrm{Tr}(\rho H_{\partial F}^2)
\]
removes ambiguity associated with negative boundary-energy expectation values and gives a non-negative dynamical penalty.

\section{Scope and Limitations}
The current three-qubit analysis is an existence-oriented toy model. It does not perform a global optimization over the complete continuous factorization manifold, derive $\lambda$ microscopically, or provide an experimental validation.

\section{Conclusion}
REM formulates relational articulation as a variational competition between informational differentiation and dynamical boundary cost. The quadratic form makes this competition sign definite and provides a basis for open-system and experimental extensions.

\begin{thebibliography}{9}
\bibitem{rovelli} C. Rovelli, Relational Quantum Mechanics, Int. J. Theor. Phys. 35, 1637 (1996).
\bibitem{zurek} W. H. Zurek, Decoherence, einselection, and the quantum origins of the classical, Rev. Mod. Phys. 75, 715 (2003).
\bibitem{carroll} S. Carroll and A. Singh, Quantum Mereology, Phys. Rev. A 103, 022213 (2021).
\bibitem{rem5} Y. Maruko, Relational Emergence Model: Variational Selection of Tensor-Product Structure in Quantum Systems (2026).
\end{thebibliography}

\end{document}
